%mac
\immini{\subsubsection{ĐỊNH LUẬT 3 NEWTON}
\setcounter{paragraph}{0}
\begin{itemize}
\item Lực không tồn tại riêng lẻ. Các lực hút hoặc đẩy luôn xuất hiện thành từng cặp giữa hai vật.
\item \textit{Trong mọi trường hợp, khi vật A tác dụng lên vật B một lực thì vật B cũng tác dụng trở lại vật A một lực. Hai lực này cùng giá, cùng độ lớn, ngược chiều}
\begin{align*}
\boder{\vec F_{AB}=-\vec F_{BA}}
\end{align*}
\end{itemize}}{\begin{tikzpicture}
\node[] at (0,0){\includegraphics[width=5cm]{images/16-1}};
\end{tikzpicture}}	
\immini{\subsubsection{CÁC ĐẶC ĐIỂM CỦA LỰC VÀ PHẢN LỰC}
\begin{itemize}	
\item Trong tương tác giữa hai vật, một lực gọi là lực tác dụng còn lực kia gọi là phản lực.
\item Đặc điểm của lực và phản lực:
\begin{itemize}	
\item Lực và phản lực luôn xuất hiện thành từng cặp (xuất hiện và mất đi đồng thời).
\item Lực và phản lực cùng tác dụng theo một đường thẳng, cùng độ lớn nhưng ngược chiều (hai lực như vậy là hai lực trực đối).
\item Lực và phản lực không cân bằng nhau (vì chúng đặt vào hai vật khác nhau).
\end{itemize}
\end{itemize}}{\begin{tikzpicture}
\node[] at (0,0){\includegraphics[width=4cm]{images/16-3}};
\end{tikzpicture}}
\immini{
\begin{itemize}	
\item Cặp lực và phản lực là hai lực cùng loại.
Một ô tô chuyển động trên mặt đường, lực do ô tô tác dụng lên mặt đường có độ lớn bằng lực mà mặt đường đẩy ô tô nhưng chúng không "khử nhau".
\end{itemize}}{
\begin{tikzpicture}
\node[] at (0,0){\includegraphics[width=4.5cm]{images/16-4}};
\end{tikzpicture}
}	
